% s-330-p330-m001; original PDF 330, 331
\Aphorism{s-330-p330-m001}{380}
Aufgabe der Philosophie ist die Enträthselung der Welt: dabei wird es nun wohl bleiben müssen. Aber man kann ein Räthsel lösen oder beseitigen, in beiden Fällen hat man etwas „enträthselt“, und ehe ein neuer Oedipus vor die alte Sphinx tritt — sollte es da nicht am Platze sein, dieses fabelhafte Ungeheuer von hinten zu besehen, ob wirklich die Jungfrau in einen Löwenleib ausläuft? Muss man vielleicht Thebaner, wenigstens Böotier sein, um an ihr etwas Gefährliches, Sphinxhaftes, Räthselhaftes zu finden? — Im Ernst, wir haben uns von unseren Philosophen viel zu viel „problematische Natur“ einreden lassen, und so ehrenwerth es sein mag, wenn wir den Einzelheiten des Lebens gegenüber das Fragen nicht verlernen: als Ganzes und Allgemeines, fürchte ich, ist die Welt das glatteste und durchsichtigste aller Rechenexempel. Wir Menschen, mit unserer starken Schwerhörigkeit gegenüber Identitäten, wir Kinder, die das Küchengeräth analytischer Urtheile auf die Strasse werfen, um die Scherben von einem synthetischen Topfflicker wieder zusammenheften zu lassen, wir haben erst ein Geheimniss aus der Welt gemacht, „über das Geschick“, wie Homer sagt — ein so verwickeltes Geheimniss, dass wir vor lauter Bedenklichkeit und Räthselratherei zuletzt nicht mehr wussten wo aus noch ein, und uns in unserer Begriffsstutzigkeit allmählich bis zu einem Gott hinauf übertölpelten. In Zeiten der Freigeisterei nennt man’s nicht Gott, sondern minder verfänglich „allgemeine Gesetzmässigkeit des Weltalls“; der Thatbestand aber bleibt derselbe, nämlich dass wir eine Causalität bestaunen, wo eine blosse Identität wirkt, dass wir ein reales Räthsel lösen wollen, wo nur ein scheinbares wegzubringen ist, dass wir eine blaue Brille tragen und uns unbändig verwundern, alle Dinge blau zu sehn —— kurz dass wir einen erhabenen Scharfsinn verschwenden, um hinter das Geheimniss des Satzes zu kommen:
\par

\begin{BookVerse}
\VerseLine{0}{Es lebt kein Schurk’ im ganzen Dänemark,}
\VerseLine{0}{Der nicht ein ausgemachter Bube wäre!}
\end{BookVerse}
